\documentclass[11pt,a4paper]{article}
\usepackage[margin=24mm]{geometry}
\usepackage{graphicx,booktabs,amsmath,amssymb,microtype}
\usepackage[colorlinks=true,linkcolor=black,citecolor=black,urlcolor=blue]{hyperref}
\usepackage{xcolor}
\usepackage{listings}
\lstset{basicstyle=\ttfamily\scriptsize,breaklines=true,columns=fullflexible,
        frame=single,framesep=4pt,showstringspaces=false,
        numbers=left,numberstyle=\tiny\color{gray},xleftmargin=16pt}
\usepackage{caption}
\captionsetup{font=small,labelfont=bf}
\newcommand{\ED}{E_{2\mathrm{D}}}

\title{\textbf{Blisters, bubbles and protrusions in 2D materials}\\[2mm]
\large A working note on toy molecular dynamics, AIREBO and LAMMPS ---\\
what each model is, what it can be trusted for, and how to check it}
\author{Prepared for Peter Bøggild (DTU) \\ \small Digital Science Workstation}
\date{29 August 2026}

\begin{document}
\maketitle

\begin{abstract}
\noindent
This note covers three things that turned out to be one thing. First, what the
DSW \texttt{graphene-md} toy model actually is --- Morse bonds, angular springs
and a Lennard-Jones substrate --- and how that differs from AIREBO and from what
LAMMPS brings. Second, the mechanics of pressurised blisters in 2D membranes,
from the Hencky solution through the adhesion balance that Bunch and co-workers
used to measure the adhesion energy of graphene. Third, and most usefully, a
method for validating a fast approximate model against an expensive accurate
one, including the ways such a validation can look convincing while proving
nothing. All measurements quoted are from runs performed for this note.
\end{abstract}

\tableofcontents

\section{The two models}

\subsection{The toy model: Morse plus angular springs plus Lennard-Jones}

Three ingredients, each a pair or triplet interaction with an explicit formula.

\paragraph{Morse bonds.} Nearest neighbours interact through
\begin{equation}
U_\mathrm{M}(d) = D_e\left[\left(1-e^{-\alpha(d-r_e)}\right)^2 - 1\right],
\end{equation}
a spring near $r_e$ that softens and goes flat at large separation. Because it
flattens, it can represent a bond \emph{breaking}: past a chosen extension the
restoring force falls away. The neighbour list is fixed at build time, so the
model knows which pairs are ``bonded'' and never discovers new ones except
through an explicit re-forming rule.

\paragraph{Angular springs.} A honeycomb held together only by central forces
between nearest neighbours has \emph{no shear rigidity} --- it is a mechanism,
not a solid, and it collapses. The toy model restores stiffness with harmonic
springs to second neighbours at $\sqrt{3}\,r_e$ (equivalently, a bond-angle
term). This is a device, not chemistry: it produces a stable lattice with
roughly the right elastic constants and nothing more.

\paragraph{Lennard-Jones substrate.} The sheet is attracted to a rigid
substrate by a 12--6 potential, truncated at $3\sigma$. This provides adhesion,
and it is the term that decides whether a blister stays stuck down.

\paragraph{Bending.} A crude ``umbrella'' term penalising out-of-plane
displacement of an atom relative to its neighbours. It has no LAMMPS analogue
and is the one term the plugin's LAMMPS exporter cannot translate.

\subsection{AIREBO: a bond-order potential}

AIREBO is a different kind of object. Its core is REBO, a \emph{bond-order}
potential in the Tersoff--Brenner family:
\begin{equation}
U = \sum_{i}\sum_{j>i}\Big[ V^{R}(r_{ij}) - b_{ij}\,V^{A}(r_{ij}) \Big],
\end{equation}
a repulsive term minus an attractive term whose strength $b_{ij}$ depends on
the \emph{local environment} --- how many neighbours atom $i$ has, at what
angles, and how their own bonding looks. That single feature is the difference
that matters:

\begin{itemize}
\item \textbf{There is no bond list.} Bonding is recomputed from geometry every
      step. Atoms form and break bonds because the environment changed, not
      because a distance crossed a threshold.
\item \textbf{Coordination is energetic.} A three-fold carbon and a two-fold
      carbon are genuinely different, so edges, vacancies and reconstructions
      come out of the potential rather than being imposed.
\item \textbf{Angles are implicit.} No separate angular spring is needed; the
      bond order carries the angular dependence.
\end{itemize}

AIREBO adds to REBO an \emph{adaptive} Lennard-Jones term (switched off between
bonded neighbours, on between non-bonded ones, so interlayer vdW works without
corrupting covalent bonding) and a torsional term for dihedral angles. Hence
\textbf{A}daptive \textbf{I}ntermolecular \textbf{REBO}.

The price is cost: the bond order requires neighbour-of-neighbour sums, so
AIREBO is one to two orders of magnitude more expensive per atom-step than a
pair potential with a fixed bond list.

\subsection{What each is good for}

\begin{table}[h]\centering\small
\begin{tabular}{lll}
\toprule
Question & Toy model & AIREBO \\
\midrule
membrane stretching, blister shape & \textbf{reliable} & reliable \\
adhesion-driven draping, wrinkling & reliable & reliable \\
phonons, elastic constants & reliable if fitted & reliable \\
bending near a peel front & crude & reliable \\
bond breaking, tearing, edges & \textbf{a distance rule, not chemistry} & reliable \\
defects, reconstruction, chemistry & \textbf{no} & reliable \\
speed & \textbf{25 steps/s at 233k atoms} & 9 steps/s \\
\bottomrule
\end{tabular}
\caption{The split is not ``approximate vs exact'' but ``mechanics vs
chemistry''. Where the physics is membrane mechanics, the cheap model is not an
approximation to the expensive one --- it contains the same physics.}
\end{table}

\section{What LAMMPS adds over a toy MD}

A purpose-built simulation does one thing quickly. LAMMPS is a framework, and
the difference shows up in five places.

\begin{enumerate}
\item \textbf{A potentials library.} AIREBO, ExTeP, REBO-MoS$_2$, Stillinger--Weber,
      ReaxFF, machine-learned potentials --- all published, parameterised and
      tested. Writing one yourself is a research project.
\item \textbf{Fixes as composable physics.} Thermostats and barostats, real
      ensembles (NVE/NVT/NPT), constraints, external forces, moving boundaries,
      bond breaking and creation. Each is a well-tested module you switch on.
\item \textbf{Neighbour machinery.} Binning, skins, rebuild criteria, exclusion
      lists, and the diagnostics to see when it goes wrong.
\item \textbf{Scale.} MPI domain decomposition to thousands of cores; the toy
      model is one shared-memory node.
\item \textbf{Analysis.} Computes, ensemble averages, stress per atom,
      correlation functions, dumps in formats other tools read.
\end{enumerate}

What it costs is generality overhead, and it is measurable. On the identical
Morse system LAMMPS spent \textbf{43\,\% of its time in \texttt{Modify}} ---
\texttt{fix move variable} evaluating an \texttt{exp()} per substrate atom per
step --- against 38\,\% in \texttt{Pair}. On the identical AIREBO system,
\textbf{Pair rose to 73\,\%} and Modify fell to 23\,\%. The framework overhead
is a fixed cost that stops mattering once the potential is genuinely expensive.
Correspondingly the plugin's speed advantage falls from $1.9\times$ on the toy
model to $1.15\times$ on AIREBO.

\section{Blister mechanics}

\subsection{The Hencky solution}

A clamped circular membrane of radius $a$ under uniform pressure $p$, with 2D
modulus $\ED$, deflects at its centre by $\delta$ where
\begin{equation}
\frac{p\,a}{\ED} = K(\nu)\left(\frac{\delta}{a}\right)^{3},
\qquad K \approx 3.09 \ \ (\nu = 0.16),
\label{eq:hencky}
\end{equation}
with shape $w(\rho) = \delta\,(1-\rho^2)^{2/3}$, $\rho = r/a$, and enclosed
volume
\begin{equation}
V = C\,\pi a^2 \delta, \qquad C \approx 0.52 .
\end{equation}
The cubic in $\delta$ is the signature of a \emph{membrane}: stiffness comes
from stretching, not bending, so the structure gets stiffer as it deflects.

\subsection{The adhesion balance}

Blisters are not clamped --- they are stuck down, and they can peel. Integrating
$p\,\mathrm{d}V$ at fixed $a$ using \eqref{eq:hencky} gives the stored strain
energy as exactly one quarter of the pressure--volume work, so the potential
energy is $\Pi = -\tfrac34 pV$. The energy release rate at the delamination
front follows from $G = -(2\pi a)^{-1}\,\partial\Pi/\partial a$:
\begin{equation}
\boxed{\;G = \tfrac54\,C\,p\,\delta \;\approx\; 0.65\,p\,\delta\;}
\label{eq:G}
\end{equation}
Setting $G = \Gamma$, the adhesion energy per unit area, and eliminating $p$
with \eqref{eq:hencky}:
\begin{equation}
\boxed{\;\frac{\delta}{a} = \left[\frac{4\Gamma}{5\,C\,K\,\ED}\right]^{1/4}
   \approx \left(\frac{\Gamma}{680\ \mathrm{N/m}}\right)^{1/4}\;}
\label{eq:aspect}
\end{equation}
and, at that aspect ratio,
\begin{equation}
\boxed{\;p \;\propto\; \frac{\Gamma^{3/4}\ED^{1/4}}{a}\;}
\end{equation}

Two consequences worth internalising. The aspect ratio is
\textbf{size-independent} --- no length appears in \eqref{eq:aspect} --- which is
why measured graphene bubbles cluster around $\delta/a \sim 0.1$--$0.2$ across
orders of magnitude in radius. And it depends on adhesion only as a
\textbf{fourth root}, so a blister is a blunt instrument for measuring $\Gamma$:
a factor of two error in $\delta/a$ is a factor of sixteen in $\Gamma$.

This is the basis of the constant-$N$ blister test of Koenig, Boddeti, Dunn and
Bunch, who used it to obtain $\Gamma = 0.45\,\mathrm{J/m^2}$ for monolayer
graphene on SiO$_2$ and $0.31\,\mathrm{J/m^2}$ for two-to-five layers.

\subsection{Constant pressure versus constant $N$}

If the pressure is held fixed, inflation is unstable: the gas does unlimited
work. If instead a \emph{fixed number of molecules} is trapped, then
$p = Nk_BT/V$, and inflating raises $V$ and drops $p$. That negative feedback is
what makes a blister an equilibrium object rather than a runaway.

This matters for simulation, and it was the origin of a bug in ours: the
protrusion driver applied a \emph{prescribed} pressure schedule, so the blister
inflated without limit and was still rising when the schedule had returned the
pressure to zero. Replacing it with a fixed-$N$ gas law, measuring $V$ from the
sheet each step, made the pressure fall from 2492 to 557\,MPa as the blister
grew --- the feedback appearing where before there was none.

\subsection{How the pressure is actually implemented, and how far to trust it}
\label{sec:gasreal}

The gas is not simulated. It is a boundary condition, applied as follows: at
each step the pocket volume is measured from the sheet,
\begin{equation}
V = \sum_{i \,:\, r_i < R} \max\!\left(z_i - z_0,\, 0\right) A_\mathrm{atom},
\end{equation}
the pressure is taken from the isothermal ideal gas law $p = Nk_BT/V$ at fixed
$T = 300$\,K, and a \emph{uniform vertical} force $pA_\mathrm{atom}$ is added to
every sheet atom inside a hard-edged disc of radius $R$. $N$ is fixed once, from
the requested pressure at the Hencky equilibrium volume.

Six approximations are buried in that paragraph, in decreasing order of how much
they should worry you.

\paragraph{1. The ideal gas law is unphysical at these pressures.} At 1\,GPa and
300\,K, $n = p/k_BT = 241\,\mathrm{nm^{-3}}$ --- roughly \emph{fourteen times}
the number density of liquid nitrogen. No real gas can be that dense; the
molecules do not fit. So the ``$N$ molecules'' in the model are a bookkeeping
device, not a physical count.

What survives is the \emph{stiffness law}: the model asserts
$p \propto V^{-1}$, i.e. $\mathrm{d}p/\mathrm{d}V = -p/V$. A real fluid at
liquid-like density is far stiffer than that --- excluded volume makes $p$ rise
much faster as $V$ shrinks. Our blister is therefore \textbf{too soft against
volume change}. Practically: a real gas would resist the initial inflation more,
snap through less violently, and ring down differently. A van der Waals form,
$p = Nk_BT/(V-Nb) - a(N/V)^2$, would cost one line and remove most of the error.

\paragraph{2. The footprint is fixed and cannot follow a peel front --- and
peeling is the whole mechanism.} The gas pushes only inside $r < R$, always. This
is not one limitation among several: \textbf{the blister test measures $\Gamma$
\emph{because} the contact line moves}. Equation~\eqref{eq:G} comes from
$\partial \Pi/\partial a$ at a free delamination front, so the radius must be an
\emph{output} of the calculation. Pinning it is differentiating with respect to a
variable the model does not allow to vary.

Two things save the result in practice, and they are worth separating.

First, \textbf{the peel is present, just not where the gas is}. The sheet outside
$R$ is genuinely free: lifted by membrane tension, held down by adhesion, with a
contact line that finds its own radius. That is why the blister behaves as though
its radius were 22\,\% larger than the footprint. What is missing is not the
moving contact line but the gas loading over the newly opened annulus.

Second, \textbf{we can check the answer, which an experiment cannot}. Koenig
\textit{et al.}\ had to trust the blister analysis because $\Gamma$ was the
unknown quantity. Here $\Gamma$ is an \emph{input}: integrating the 12--6
substrate potential over a continuum sheet gives
\begin{equation}
\Gamma_\mathrm{LJ} = 1.2\,\pi\, n_t n_s\, \varepsilon\, \sigma^2
 = 0.246\ \mathrm{J/m^2} \quad (0.241 \ \text{with the } 3\sigma \text{ cutoff}),
\end{equation}
at an equilibrium separation of exactly $d = \sigma = 3.42$\,\AA. The blister
test returned $0.22\,\mathrm{J/m^2}$ --- \textbf{9\,\% low}, with the sign one
would predict: a pinned footprint pressurises less area than a true blister of
the same delaminated radius, so $p\delta$ slightly under-counts $G$.

That the cutoff costs only 2.1\,\% of the adhesion energy also disposes of a
plausible-sounding worry: extending the Lennard-Jones range would not
meaningfully change any threshold in this note.

\paragraph{The fix.} Replace the $r<R$ test with a \emph{gap} test: pressurise
wherever the sheet stands more than about a molecular diameter above contact, and
count that volume. The gas then occupies the void it actually occupies, the
pressurised region follows the delamination front automatically, and the blister
radius becomes an output. A small pre-lifted patch is needed to nucleate, since a
perfectly adhered sheet encloses no volume. With that change the model would be a
blister test by construction rather than by accident, and the 9\,\% could be
re-measured to see whether it closes.

\paragraph{3. Volume is counted only inside the footprint.} Any lift outside $R$
contributes gas volume in reality but not in the model, so $p$ is
correspondingly overestimated. Consistent with (2), and in the same direction.

\paragraph{4. The force is vertical, not normal to the surface.} Real pressure
acts along the local surface normal. At $\delta/a \approx 0.13$ the steepest
part of the profile is inclined by roughly $20$--$25^\circ$, so near the rim the
model over-counts the vertical component by $\sim$10\,\% and omits an inward
radial component of comparable size. The rim is exactly where the adhesion
balance is decided, so this is not as harmless as its magnitude suggests.

\paragraph{5. Isothermal, with no thermal coupling.} $T$ is pinned at 300\,K.
A real snap-through expansion is closer to adiabatic: the gas would cool, $p$
would fall further than $V^{-1}$ alone predicts, and the overshoot would be
smaller. Our ring-down is more violent than reality.

\paragraph{6. No leakage --- and here the model is right.} Fixed $N$ assumes a
perfect seal. For graphene that is a genuinely good idealisation: Bunch
\textit{et al.} showed graphene membranes to be impermeable to standard gases,
which is what makes the constant-$N$ blister test possible experimentally in the
first place.

\paragraph{What this does \emph{not} contaminate.} The adhesion energy extracted
in Section~\ref{sec:gamma} is safe from all of the above, because
$G = \tfrac54 C p \delta$ uses the pressure and deflection \emph{measured at
equilibrium}. Whatever route the gas took to get there, and whatever equation of
state generated that pressure, the balance at equilibrium involves only $p$ and
$\delta$. The gas law determines \emph{which} equilibrium you land on; it does
not change the condition that defines it. Results that depend on the path ---
the snap-through amplitude, the ring-down, the dynamic threshold pressure ---
are a different matter and should be treated as model behaviour.

\section{What we measured}

\subsection{The exporter reproduces the plugin}

The strongest available check on two implementations is to put them in the
\emph{same state} and compare an extensive quantity. Exporting a LAMMPS deck
from the plugin and evaluating step 0:

\begin{table}[h]\centering\small
\begin{tabular}{lrrr}
\toprule
System & Plugin $E_\mathrm{pot}$ & Stock LAMMPS & rel.\ difference \\
\midrule
930\,629 atoms, Morse  & $-14591.6$ eV & $-14591.569$ eV & $2\times10^{-6}$ \\
232\,735 atoms, Morse  & $-3655.11$ eV & $-3655.106$ eV  & $1\times10^{-6}$ \\
232\,735 atoms, AIREBO & $-711418$ eV  & $-711418.41$ eV & $6\times10^{-7}$ \\
\bottomrule
\end{tabular}
\end{table}

AIREBO's cohesive energy comes out at $-7.44$\,eV/atom against the published
$-7.408$ for infinite graphene, the balance being substrate adhesion --- an
independent check that the potential is correctly parameterised.

\subsection{Blister geometry across engines}

For the twisted-bubble case the four engines split cleanly:

\begin{itemize}
\item \textbf{In-plane} quantities are engine-insensitive. Lattice registry
      statistics agreed to four decimal places (mean $0.6668$, s.d.\ $0.2717$)
      between the plugin's embedded AIREBO and stock LAMMPS, and the atoms
      differed by at most $0.097$\,\AA{} in plane.
\item \textbf{Out-of-plane} quantities are not. The same two runs differed by
      $1.58$\,\AA{} vertically, and blister heights differed by up to 15\,\%.
\end{itemize}

\subsection{A stable blister, and $\Gamma$}
\label{sec:gamma}

With the fixed-$N$ gas, a stable blister exists in a narrow window. At
$R = 3$\,nm and 700\,MPa the volume drift fell to 1.7\,\% over 49\,000 steps,
settling at $\delta = 5.13$\,\AA, $\delta/a = 0.171$, at a settled pressure of
658\,MPa. Applying \eqref{eq:G}:
\begin{equation}
\Gamma = 0.65\,p\,\delta
       = 0.65 \times 658\times10^{6} \times 5.13\times10^{-10}
       = 0.22\ \mathrm{J/m^2},
\end{equation}
The number to compare against is not Koenig's experimental value but the
adhesion energy we \emph{put into} the model,
$\Gamma_\mathrm{LJ} = 0.241\,\mathrm{J/m^2}$ (Section~\ref{sec:gasreal}): the
blister test recovers its own input to \textbf{9\,\%}, and low, as a pinned
footprint should make it. It also happens to sit near the
$0.24\,\mathrm{J/m^2}$ Koenig \textit{et al.}\ report, but that is a different
interface and the closeness is a coincidence of parameterisation, not evidence.

Feeding that $\Gamma$ back through \eqref{eq:aspect} predicts runaway near
250\,MPa at $R = 10$\,nm; the measured transition was 200--260\,MPa.

\subsection{The rule tested at a different size}

The strongest test of \eqref{eq:aspect} is prediction rather than fitting. Taking
$\Gamma = 0.22\,\mathrm{J/m^2}$ from the $R = 3$\,nm blister above, the rule
predicts $\delta/a = 0.134$ at \emph{any} radius. Running $R = 10$\,nm at
160\,MPa --- a different size, and a quarter of the pressure --- gives a settled
$\delta/a = 0.128$ (Figure~\ref{fig:blister10}), within 4\,\% of the prediction.

The approach to equilibrium is itself instructive. The sheet stays pinned flat
while the gas fills, then delaminates in a \emph{snap-through} at around step
10\,000, overshoots to 24\,\AA, and rings down to $12.84 \pm 1.8$\,\AA{} over
the following 60\,000 steps, the oscillation amplitude decaying by a factor of
fourteen. The snap is physical: adhesion pins the membrane until the gas load
exceeds the peel threshold, and then releases abruptly.

\begin{figure}[h]\centering
\includegraphics[width=\textwidth]{figs/blister10-settled.png}
\caption{A settled fixed-$N$ blister, $R = 10$\,nm at 160\,MPa. Left: plan view,
dashed circle marking the gas footprint. Centre: the measured radial profile
against the Hencky shape through the same peak --- it tracks near the centre and
falls below from $r \approx 40$\,\AA{} outward, as a peeling rim should. Right:
centre height against step, showing the snap-through and the ring-down.}
\label{fig:blister10}
\end{figure}

\paragraph{A measurement artefact worth naming.} The automated settling test
reported ``still moving'' with 13\,\% drift, because it compared the means of
two successive windows --- a statistic that a \emph{decaying oscillation}
defeats. The mean was in fact stationary (a linear fit over the last half gives
$-0.55$\,\AA{} across 35\,000 steps) while the amplitude fell from 6.05 to
0.43\,\AA. The system had settled; the test could not see it. This belongs with
the failure modes of Section~5.2.

\paragraph{The catch.} Extracting $\Gamma$ the \emph{other} way --- from the
measured aspect ratio through \eqref{eq:aspect} --- gives $0.58\,\mathrm{J/m^2}$,
a factor $2.6$ different. The two routes disagree because our blister does not
satisfy \eqref{eq:hencky}: it sits 31\,\% above the Hencky height at $R = 3$\,nm.
The Koenig extraction assumes the membrane obeys Hencky. When it does not, the
choice of extraction route silently changes the answer threefold, and nothing in
either calculation warns you.

The discrepancy is a \emph{boundary condition}, not a modulus. Hencky clamps the
rim; a real blister peels, so the effective radius exceeds the gas footprint, and
since $\delta \propto a^{4/3}$ a 22\,\% larger effective radius produces exactly
the observed 31\,\% excess. Two pieces of evidence support this over the obvious
alternative that the toy model is too soft. First, running the identical case on
AIREBO --- which has the correct graphene stiffness --- gave $5.105$\,\AA{}
against the toy model's $5.131$\,\AA, agreeing to 0.5\,\%. Second, the excess
\emph{shrinks with size}: $+31\,\%$ at $R = 3$\,nm but $+8.3\,\%$ at
$R = 10$\,nm, exactly as a rim correction should.

\section{Twisted bubbles: registry and strain}

\subsection{A single sheet, with a free peel front}

Replacing the fixed footprint with one that follows the peel front
(Section~\ref{sec:gasreal}) makes the delamination radius an output. For a
$2^\circ$ twisted sheet with a 4\,nm seed at 600\,MPa the front settled at
45\,\AA{} against an adhesion-limited prediction of 42.6\,\AA.

The distortion is not where one might guess. The bubble interior is in mild
tension ($\mathrm{tr}\,\varepsilon \approx +0.017$ at the centre), but the
dominant features are a \textbf{ring of compression reaching $-0.045$ and a ring
of shear reaching $0.031$ sitting exactly at the delamination edge}, with the
adhered region beyond essentially unstrained. That is the classic blister
signature --- radial tension, azimuthal compression at the rim --- resolved atom
by atom.

The moiré is also visibly smeared inside the bubble. Registry depends only on
in-plane position, so that is real lateral displacement: the sheet stretches
radially to accommodate the bulge and slides out of registry with the substrate.

\subsection{A genuine bilayer, and where the gas sits}

With two live layers the placement of the gas changes the answer completely.
Identical twist, seed and pressure; only \texttt{gasWhere} differs.

\begin{table}[h]\centering\small
\begin{tabular}{lrrrr}
\toprule
& \multicolumn{2}{c}{peak height (\AA)} & \multicolumn{2}{c}{dilatation range} \\
\cmidrule(lr){2-3}\cmidrule(lr){4-5}
Placement & layer 1 & layer 2 & layer 1 & layer 2 \\
\midrule
gas \textbf{between} the layers & 3.53 & 12.35 & $-0.003 \dots +0.001$ & $-0.023 \dots +0.006$ \\
gas \textbf{below} the bilayer  & 7.54 & 10.94 & $-0.011 \dots +0.003$ & $-0.011 \dots +0.003$ \\
\bottomrule
\end{tabular}
\caption{Gas between the layers puts essentially all the deformation into the
upper layer --- the lower one stays flat and ten times less strained. Gas below
deforms both almost identically.}
\end{table}

Two consequences follow. First, the below case is \textbf{1.4\,\AA{} lower at
the same pressure} (10.94 against 12.35\,\AA), because two layers are being bent
rather than one: the effective bending rigidity is higher. Second, the
interlayer gap behaves quite differently:

\begin{table}[h]\centering\small
\begin{tabular}{lrr}
\toprule
& gap at the bubble centre & gap in the far field \\
\midrule
gas between, 600\,MPa  & 8.28\,\AA  & 3.42\,\AA \\
gas between, 1400\,MPa & 13.04\,\AA & 3.42\,\AA \\
gas below, 600\,MPa    & 3.40\,\AA  & 3.42\,\AA \\
gas below, 1400\,MPa   & 3.41\,\AA  & 3.42\,\AA \\
\bottomrule
\end{tabular}
\caption{A bubble between the layers prises them apart, so $h \gg 3.35$\,\AA{}
over the bubble is physics, not an artefact. A bubble below carries the pair up
together and the spacing never departs from equilibrium, even with the bilayer
lifted 10.7\,\AA{} off the substrate. The far field is 3.42\,\AA{} in every
run, which is the built-in 3.35\,\AA{} plus thermal roughness.}
\end{table}

\begin{figure}[h]\centering
\includegraphics[width=\textwidth]{figs/bilayer-layers.png}
\caption{Dilatation in both layers for each placement, with the layer height
profiles. Top row, gas between: the strain is confined to layer 2. Bottom row,
gas below: the two layers carry nearly identical strain and their height
profiles are parallel curves a constant 3.35\,\AA{} apart.}
\end{figure}

\begin{figure}[h]\centering
\includegraphics[width=\textwidth]{figs/bilayer-compare.png}
\caption{Height, registry, dilatation and shear for the driven layer in each
placement. The strain rings sit at the seed radius in every panel.}
\end{figure}

\paragraph{A threshold that differs from the monolayer.} At 600\,MPa the
contact line stayed pinned at the seed in both bilayer runs, where a single
sheet at the same seed and pressure advanced from 40 to 45\,\AA. Raising the
pressure to 1400\,MPa moved the bilayer front out to 62\,\AA. The plausible
reading is that a bilayer's lower layer can dish downward and absorb part of the
gas work, so less of it goes into opening new interface --- but that is an
interpretation of two pressures, not a measured threshold, and it deserves a
proper sweep before being believed.

\section{How to validate a toy model against LAMMPS}

\subsection{The tests, strongest first}

\begin{enumerate}
\item \textbf{Same state, same energy.} Export the exact configuration and
      compare total potential energy. This is by far the strongest test: it is
      a single number summing every interaction, and agreement to six figures
      is very hard to obtain by accident. Everything else is weaker.
\item \textbf{Known physical constants.} Cohesive energy per atom, lattice
      constant, elastic moduli, phonon frequencies. These catch a
      mis-parameterised potential that is nonetheless self-consistent.
\item \textbf{An analytic limit.} Hencky, a harmonic oscillator, an isolated
      dimer --- anything with a closed-form answer.
\item \textbf{Observable geometry.} Heights, profiles, aspect ratios. Weaker,
      because many wrong models produce plausible shapes.
\item \textbf{Trajectory agreement.} Effectively worthless. MD is chaotic;
      correct implementations diverge atom-for-atom within hundreds of steps,
      and float32 versus float64 is enough to do it.
\end{enumerate}

\subsection{Five ways a validation lies}

Each of these occurred while preparing this note.

\paragraph{The vacuous test.} A check that ``no surviving bond exceeds the break
threshold'' passed with a maximum strain of $0.0000$ --- because every bond had
broken and the set was empty. \emph{Always assert that the set under test is
non-empty}, and prefer conditions that would fail if the system were destroyed.

\paragraph{The unconverged measurement.} A pressure scan reported an aspect
ratio of $0.083$ at $R = 10$\,nm, suggesting the size-independence of
\eqref{eq:aspect} was violated. The runs were time-limited to $\sim$4700 steps.
Given ten times as long the same case reached $0.129$ and was still rising
toward the predicted $0.134$. \emph{A quantity that has not stopped changing is
not a measurement.}

\paragraph{The shared term.} Toy model and AIREBO agreed on blister height to
0.5\,\%, which looks like strong mutual validation. But the two engines differ
\emph{only} in the in-plane potential: they share the same Lennard-Jones
substrate adhesion and the same gas model, both of which live in the plugin
rather than in the potential. So the agreement validates the membrane half and
says nothing about adhesion --- and adhesion is precisely what sets the boundary
condition responsible for the 31\,\% Hencky discrepancy. \emph{Two models agree
most easily about the parts they share.}

\paragraph{The statistic that cannot see the answer.} A settling test built on
the difference between two window means reported 13\,\% ``drift'' for a system
whose mean was flat to $0.55$\,\AA{} over 35\,000 steps. It was measuring a
decaying oscillation, which that statistic cannot distinguish from a trend.
\emph{Choose a convergence statistic that matches how the system actually
approaches equilibrium} --- for a ringing system, fit the trend and track the
amplitude separately.

\paragraph{The proxy that stops meaning what it meant.} A stability criterion
written as ``the blister centre is below the Lennard-Jones cutoff'' was sound
for a 3\,nm blister and nonsense for a 10\,nm one, where a correct blister is
necessarily taller than the cutoff. It silently relabelled every large blister
as ``escaped''. \emph{A threshold derived for one regime must be re-derived, not
rescaled, for another.}

\subsection{A working procedure}

\begin{enumerate}
\item Build the configuration once, in the fast model.
\item Export it and confirm the energy matches to five or six figures. If it
      does not, stop --- nothing downstream is meaningful.
\item Identify which terms the two codes \emph{share}, and state explicitly that
      those are untested by the comparison.
\item Compare a physical observable, with a convergence criterion measured, not
      assumed.
\item Check against an analytic limit where one exists, and when it disagrees,
      determine whether the model or the analytic assumption is at fault.
\item Use the fast model for exploration and the expensive one for spot checks
      at the settings that matter.
\end{enumerate}

\section{Practical notes}

Tooling for the above lives in \texttt{DSW/tools/}: \texttt{cyclerun.js} (build,
drive, export, capture), \texttt{gassweep.js} and \texttt{gasconfirm.js} (blister
equilibria and Hencky comparison), \texttt{registry.py} (registry computed
identically for plugin and LAMMPS data), \texttt{makemovies.py} and
\texttt{blister\_movie.py}. \texttt{DSW/tools/README.md} lists the traps that
cost real time.

For simulating a membrane in a material with no classical potential ---
$\alpha$-RuCl$_3$ being the case at hand --- the ordering follows from
Section~1.3. Blister mechanics needs only $\ED$, bending rigidity and $\Gamma$,
so refitting the toy model to those three numbers gives a usable membrane
quickly and will be honestly wrong about chemistry. A Stillinger--Weber fit to
DFT phonons is the middle route, and the three-type reduced realisation already
used for WS$_2$, MoTe$_2$ and WTe$_2$ applies unchanged to a Cl--Ru--Cl
trilayer. A universal machine-learned potential is the best fidelity available.
None of them is magnetic, which for a Kitaev material is a real limitation.

\appendix

\section{A worked LAMMPS deck, line by line}
\label{sec:deck}

This is not an illustrative script. It is the deck the plugin exported for a
\textbf{twisted bubble} --- $2^\circ$ twist, 7.5\,nm Hencky gas pocket at
100\,kPa, 232\,497 atoms --- and it is the one that was actually run:
7\,min 06\,s on 18 threads, 226 dump frames, no errors. Its first line of output
is the check that matters:

\begin{center}\small
\begin{tabular}{lr}
\toprule
plugin $E_\mathrm{pot}$ at the moment of export & $-709090$ eV \\
LAMMPS \texttt{PotEng} at step 0 of this deck   & $-709089.85$ eV \\
\bottomrule
\end{tabular}
\end{center}

\noindent Six significant figures, from two independent codes evaluating the
same configuration. If any line below were wrong, that number would not match.

\lstinputlisting[caption={\texttt{run.in} as exported, verbatim.}]{appendix-run.in}

\subsection*{What each part does}

\paragraph{Lines 5--11, threads.} LAMMPS silently falls back to \emph{one}
OpenMP thread when \texttt{OMP\_NUM\_THREADS} is unset, which quietly wastes an
18-core machine. \texttt{suffix omp} then selects the \texttt{/omp} variants of
every style that has one.

\paragraph{Line 13, \texttt{units metal}.} Energies in eV, distances in
\AA{}, time in ps, so \texttt{timestep 0.001} is 1\,fs. Every number below is in
those units.

\paragraph{Line 15, \texttt{boundary m m m}.} Non-periodic with a
\emph{minimum} box rather than shrink-wrapped (\texttt{s}). This matters: on a
flat sheet, shrink-wrapping collapses the box in $z$ to near-zero thickness and
LAMMPS refuses to bin the neighbour list.

\paragraph{Line 16, \texttt{atom\_style atomic}.} No bond list at all. AIREBO
determines bonding from geometry every step, so there is nothing to declare ---
the difference from the classic engine's \texttt{molecular} style, which does
carry an explicit bond and angle list.

\paragraph{Lines 20--23, the potential.}
\begin{itemize}
\item \texttt{airebo 3.0} --- the 3.0 is the LJ cutoff in units of $\sigma$;
      the adaptive LJ and torsion terms are on by default.
\item \texttt{pair\_coeff * * airebo CH.airebo C NULL} --- one element per atom
      type. Type 1 is carbon; type 2 is \texttt{NULL}, i.e. the substrate takes
      no part in AIREBO at all.
\item \texttt{pair\_coeff 1 2 lj/cut 0.002387 3.42 10.26} --- the sheet's
      adhesion to the substrate, $\varepsilon = 2.387$\,meV, $\sigma = 3.42$\,\AA,
      cut at $3\sigma$. This replaces the \texttt{fix external} callback the
      plugin uses internally.
\item \texttt{pair\_coeff 2 2 none} --- the substrate does not interact with
      itself. It is rigid, so that energy is a constant and computing it would
      only cost time.
\end{itemize}

\paragraph{Lines 29--31, groups.} \texttt{sheet} and \texttt{sub} by type. The
\texttt{edge} group is empty here because the edge-spring collar is switched
off; it is emitted as \texttt{group edge empty} rather than as a list of atom
ids, because \texttt{group \dots id} rescans every atom once per argument and a
37\,000-id list once kept LAMMPS in setup for over twelve minutes.

\paragraph{Lines 33--35, integration.} \textbf{Only the sheet is integrated}
--- there is no \texttt{fix nve} on \texttt{sub}, which is exactly how a rigid
substrate is expressed. \texttt{fix viscous} supplies the plugin's damping.
Note this is \emph{not} a thermostat: it removes energy without supplying
fluctuations, so the run relaxes toward $T=0$ rather than sampling an ensemble.

\paragraph{Lines 37--43, the protrusion driver.} The interesting part, and the
one worth checking arithmetic on.
\begin{itemize}
\item \texttt{v\_tps} is elapsed time in ps.
\item \texttt{v\_elev} is a ramp--hold--return written as boolean algebra,
      because LAMMPS equal-style variables have no \texttt{if}. The three terms
      are (rise)$\times$($t<t_1$) $+$ (hold)$\times$($t_1\le t<t_2$) $+$
      (fall)$\times$($t_2\le t<t_3$). With a rise rate of 10\,\AA/ps (the
      plugin's 0.01\,\AA/step over a 1\,fs step) from 0.05\,\AA{} to the Hencky
      height 6.922\,\AA{}: $t_1 = 0.6872$\,ps, a 0.5\,ps hold gives
      $t_2 = 1.1872$, and the fall gives $t_3 = 1.8794$\,ps. That is 1879 steps,
      which is exactly the cycle length quoted at line 60.
\item \texttt{v\_gasF} is the force per atom at full pressure,
      $0.0408818$\,eV/\AA. Check it:
      $10^{8}\,\mathrm{Pa} \times 6.2415\times10^{-12} \times 2.62\,\text{\AA}^2
      \times 25 = 0.04088$. \textbf{The 25 is \texttt{betweenBoost}}, a
      deliberate fudge in the open-loop model without which a 1\,bar pocket
      does nothing at this scale. It is not physics, and the fixed-$N$ model in
      Section~\ref{sec:gasreal} does not use it.
\item \texttt{v\_wfoot} is the Hencky footprint $(1-r^2/R^2)^{2/3}$ with
      $R^2 = 5625\,\text{\AA}^2$, i.e. $R = 75$\,\AA. The nested
      $((r^2<R^2))$ inside the base is not redundant: outside the blister
      $1-r^2/R^2$ is negative, a negative base to a fractional power is NaN, and
      $0\times\mathrm{NaN}=\mathrm{NaN}$. Folding the guard into the base keeps
      it at exactly 1 outside. An earlier version guarded only the result and
      NaN'd the entire sheet on the first step.
\item \texttt{fix gas sheet addforce} applies it, upward only.
\end{itemize}

\paragraph{Lines 45--48, thermodynamics.} \texttt{compute Tsheet sheet temp} is
not cosmetic: thermo's default temperature averages over the frozen substrate
as well, and with roughly 60\,\% of the atoms frozen it reads about half the
true sheet temperature.

\paragraph{Lines 49--54, output.} The sheet is dumped ten times more often than
the substrate. Dumping every atom at the sheet's cadence would be gigabytes of
text for no extra information, since the substrate is rigid.

\subsection*{What this deck does \emph{not} reproduce}

Worth stating plainly, because it is what a reader would otherwise assume.

\begin{enumerate}
\item \textbf{This is the open-loop gas model.} Pressure follows the schedule
      regardless of the bubble's volume, so it inflates without the
      $p = NkT/V$ feedback of Section~\ref{sec:gasreal}, and the
      \texttt{betweenBoost} factor of 25 is active. It faithfully reproduces
      what the plugin did in this run --- it is not the better model.
\item \textbf{The footprint is a fixed disc.} The free peel front is
      \emph{not} exported: the exporter emits the fixed-radius form, so a deck
      cannot currently reproduce a moving contact line. Since $\Gamma$ is defined
      at a moving front, no adhesion energy should be extracted from this deck.
\item \textbf{No bilayer.} This is one mobile sheet on a rigid substrate. The
      two-live-layer runs come from a different plugin (\texttt{moire-bubble})
      which has no exporter at all, so those results exist only inside DSW.
\item \textbf{Velocities are carried but the state is not a thermal ensemble.}
      \texttt{system.data} includes a \texttt{Velocities} section so the run
      continues from the plugin's instant, but with \texttt{fix viscous} and no
      thermostat the trajectory relaxes rather than samples.
\item \textbf{Classic-engine decks omit the bending umbrella.} Not an issue
      here --- this is the AIREBO engine, where bending is carried by the
      potential --- but a deck exported from the Morse engine loses that term,
      which is why its blisters sit lower.
\end{enumerate}

\end{document}
